\documentclass[12pt]{article}
\usepackage[utf8]{inputenc}
\usepackage[T1]{fontenc}
\usepackage{hyperref}
\title{Pogolo Language Resources for K-12 Curriculum Development}
\author{Compiled by Claude Sonnet 5.5 (claude-sonnet-5-5) for Andr\'as Kornai}
\date{}
\begin{document}
\maketitle

\section{Introduction}
Pogolo (Pogoro, Chipogolo; ISO 639-3 poy) is a Bantu language of the Morogoro and Iringa regions of Tanzania. Wikipedia gives about 180,000 speakers (a 1987 figure). Little published material turned up in search.

\section{Primary Educational Resources}
No primers, textbooks or school programmes for Pogolo were found.

\section{Linguistic and Reference Materials}
\begin{itemize}
\item Overview article: \url{https://en.wikipedia.org/wiki/Pogolo_language}
\end{itemize}
No dictionary or grammar was found in the search results.

\section{Religious and Community Resources}
\begin{itemize}
\item Pogolo New Testament, Digital Bible Society: \url{https://dbs.org/bibles/POYNTP}
\item Joshua Project profile of the Pogoro people (population estimates vary between its pages): \url{https://joshuaproject.net/people_groups/14425}
\end{itemize}

\section{Digital and Audio Resources}
No audio or software resources were found.

\section{Curriculum Development Context}
The language is widely spoken but apparently not taught in schools, and the only substantial text found is the New Testament. A curriculum would need new primary materials developed with the community.
\end{document}